\documentclass{article}
\usepackage[utf8]{inputenc}
\usepackage{hyperref}
\usepackage{url}
\title{Is the PFAS-rice model usable as a dietary risk-assessment tool, and at what assurance level? A risk-assessment-readiness verdict over grain (the dietary compartment).}
\author{sci-adk (deterministic render)}
\date{\today}

\begin{document}
\maketitle

\noindent\textit{Draft compiled by sci-adk from Spec \texttt{pfas-rice-risk-readiness} (v1). Belief state is revisable as Evidence accrues.}

\begin{abstract}
This run answers whether the PFAS-rice model can be used as a dietary risk-assessment tool, and at what assurance level. Two developments make the question live. First, a breakthrough: the long-chain root-to-shoot problem was NOT structurally unresolvable -- that was an artifact of holding the xylem-loading factor f\_xy fixed and only adding the non-subtractable lipid term. A long chain has two INDEPENDENT physical properties -- a LOW f\_xy (it is strongly retained in the root) and a need for an ENHANCED active carrier (to build the high measured root uptake) -- and once f\_xy is free, the standard 2-pool reproduces C10-C12 root, straw and grain at log10 RMSE about 0.08 (a saturated, degrees-of-freedom-zero fit = structural adequacy, not a-priori prediction). The single-pool core could not even reproduce these chains. So there is no structural blind spot across the diet-relevant congener range. Second, the dietary compartment itself is predicted out-of-sample: the lipid mechanism, transferred without grain refit, predicts the independent Kim 2019 brown-rice grain BAF at log10 RMSE 0.48 (factor about 3), reliable subset 0.20 (factor about 1.6). Verdict: the model is usable as a SCREENING-level dietary risk-assessment tool, with the honest bound that the worst-case grain miss (PFOS endosperm, about 5x) keeps it screening-grade, not regulatory-precision -- dietary use must carry congener-specific uncertainty factors (about 3x typical, up to about 5x for PFSA endosperm and ether).
\end{abstract}

\section{Introduction}
The audit's standing conclusion was that the model could not support dietary risk assessment: the grain compartment was under-predicted and the long chains were unreproducible. The directive here is concrete -- the model must be usable as a risk-assessment tool -- so this run tests readiness against measured data rather than asserting it. Risk assessment for PFAS in rice is, in the first instance, about the grain (brown rice), the dietary-exposure compartment. Readiness therefore has two parts: structural coverage (can the model even represent every congener that matters, including the long chains), and predictive grain accuracy out-of-sample (does it get the dietary number right on data it was not fit to). The breakthrough supplies the first; the cross-dataset lipid result supplies the second. The engine resolves the assurance-level verdict from the frozen thresholds; the narrative is agent-authored prose input.

\section{Goal}
Is the PFAS-rice model usable as a dietary risk-assessment tool, and at what assurance level? A risk-assessment-readiness verdict over grain (the dietary compartment).

Decide, against measured data, whether the model can serve dietary risk assessment
for PFAS in brown rice:

1. Structural coverage: with the 2-pool + free f\_xy + active carrier, is there any
   remaining structural blind spot across the C4-C12 congener range (including the
   long chains)?
2. Dietary-compartment prediction: is brown-rice grain predicted OUT-OF-SAMPLE on an
   independent dataset within a screening-adequate factor (\textasciitilde{}3)?
3. Reliable subset: is the reliable-detection grain prediction within \textasciitilde{}factor 1.6?
4. Honest assurance level: is the worst-case grain residual large enough that the
   tool is SCREENING-level (bounded uncertainty), not regulatory-precision?

\section{Background}
The audit concluded the model was not yet usable for dietary risk assessment: the
risk-relevant compartment (brown-rice grain) was under-predicted (H4 REFUTED), and
the longest chains could not be reproduced (the single-pool refit hit ceilings,
\textasciitilde{}4-6x under). Two developments change this. (1) BREAKTHROUGH
(validation/longchain\_closure.py): the long-chain root<->shoot was NOT structurally
unresolvable -- that was an artifact of holding f\_xy fixed and only adding the
non-subtractable lipid term. Recognizing two INDEPENDENT physical facts about a
long chain -- a LOW xylem-loading f\_xy (strong root retention) and an ENHANCED
active carrier (high uptake) -- the standard 2-pool reproduces C10-C12 root, straw
and grain at log10 RMSE \textasciitilde{}0.08 (saturated, DOF 0 = structural adequacy, not a-priori
prediction). (2) The K\_PL-gated lipid mechanism predicts the dietary compartment
OUT-OF-SAMPLE: independent Kim 2019 brown-rice grain BAF at log10 RMSE 0.48 (factor
\textasciitilde{}3), reliable subset 0.20 (factor \textasciitilde{}1.6). This run asks whether, together, these make
the model usable as a dietary risk-assessment tool, and at what level.

\section{Method}
Run the live breakthrough experiment (validation/longchain\_closure.py) for the
long-chain structural-coverage statistic, and take the grain out-of-sample statistics
from the committed cross-dataset runs (runs/pfas-rice-oos-multidataset for Kim grain;
runs/pfas-rice-oos-lipid for the worst-case PFOS-endosperm residual). Freeze these as
threshold hypotheses; the engine resolves them and renders the paper. All evidence is
measured (vs Yamazaki / Kim / Tang). The verdict is explicitly an assurance-LEVEL
statement, not a binary -- screening vs regulatory -- so the honest worst-case bound
is a first-class hypothesis, not a footnote.

Planned approaches:
\begin{itemize}
  \item live breakthrough experiment for the long-chain structural-coverage RMSE
  \item grain out-of-sample statistics from the committed cross-dataset runs (Kim, Tang)
  \item threshold hypotheses -> engine resolves the assurance-level verdict and renders
  \item record congener-specific uncertainty factors for a risk assessor
\end{itemize}

\section{Hypotheses and findings}

\subsection{With the 2-pool + free per-congener f\_xy + active carrier, the model reproduces the full C4-C12 series INCLUDING the long chains (saturated structural-adequacy fit, log10 RMSE < 0.2) -- there is no structural blind spot across the diet-relevant congener range (the long-chain gap is closed).}
\begin{itemize}
  \item Hypothesis id: \texttt{hyp-risk-structural-coverage} (confirmatory)
  \item Decision rule (threshold): with the 2-pool + free f\_xy + active carrier, the model reproduces the full C4-C12 series including the long chains at log10 RMSE < 0.2 (no structural blind spot across the diet-relevant congener range) => support; a remaining long-chain gap => refute
  \item \textbf{Status: supported} --- confidence 0.112 (credence)
  \item Basis: threshold rule: statistic 'point'=0.081 < 0.2 is met (combine='latest', margin=0.119)
  \item \textit{Evidence validity:} referent=empirical; data\_source(s)=measured
\end{itemize}

\subsection{Brown-rice grain (the dietary compartment) is predicted OUT-OF-SAMPLE on the independent Kim 2019 dataset within a screening-adequate factor \textasciitilde{}3 (log10 RMSE < 0.5), using the lipid mechanism transferred without grain refit.}
\begin{itemize}
  \item Hypothesis id: \texttt{hyp-risk-grain-oos} (confirmatory)
  \item Decision rule (threshold): brown-rice grain (the dietary compartment) is predicted OUT-OF-SAMPLE on the independent Kim 2019 dataset within a screening-adequate factor \textasciitilde{}3 (log10 RMSE < 0.5) => support; >= 0.5 => refute
  \item \textbf{Status: supported} --- confidence 0.0198 (credence)
  \item Basis: threshold rule: statistic 'point'=0.48 < 0.5 is met (combine='latest', margin=0.02)
  \item \textit{Evidence validity:} referent=empirical; data\_source(s)=measured
\end{itemize}

\subsection{On the reliable-detection grain subset (DF>=15\%), the out-of-sample grain prediction is within \textasciitilde{}factor 1.6 (log10 RMSE < 0.3).}
\begin{itemize}
  \item Hypothesis id: \texttt{hyp-risk-grain-reliable} (confirmatory)
  \item Decision rule (threshold): on the reliable-detection grain subset (DF>=15\%), the out-of-sample prediction is within \textasciitilde{}factor 1.6 (log10 RMSE < 0.3) => support; >= 0.3 => refute
  \item \textbf{Status: supported} --- confidence 0.0952 (credence)
  \item Basis: threshold rule: statistic 'point'=0.2 < 0.3 is met (combine='latest', margin=0.1)
  \item \textit{Evidence validity:} referent=empirical; data\_source(s)=measured
\end{itemize}

\subsection{The tool is SCREENING-level, not regulatory-precision: the worst-case out-of-sample grain miss (PFOS endosperm, \textasciitilde{}5x) exceeds a factor 3 (log10 > 0.5), so dietary use must carry congener-specific uncertainty factors.}
\begin{itemize}
  \item Hypothesis id: \texttt{hyp-risk-screening-bound} (confirmatory)
  \item Decision rule (threshold): the worst-case grain residual is large enough that the tool is SCREENING-level, not regulatory-precision: the worst out-of-sample grain miss (PFOS endosperm) exceeds a factor 3 (log10 > 0.5) => support (screening-grade, bounded uncertainty); <= 0.5 => refute (would be regulatory-precision)
  \item \textbf{Status: supported} --- confidence 0.181 (credence)
  \item Basis: threshold rule: statistic 'point'=0.7 > 0.5 is met (combine='latest', margin=0.2)
  \item \textit{Evidence validity:} referent=empirical; data\_source(s)=measured
\end{itemize}

\section{Evidence}
\begin{itemize}
  \item \texttt{evi-risk-literature} (literature): finding=cited prior work: long-chain root retention / low TSCF (newcontam 2025; Adu 2024 ML 10.1021/acsestengg.4c00107); Chen 20
  \item \texttt{evi-risk-coverage} (experiment\_run): point=0.081, finding=validation/longchain\_closure.py (BREAKTHROUGH): the 2-pool + free per-congener f\_xy + active carrier reproduces the long
  \item \texttt{evi-risk-grain-oos} (experiment\_run): point=0.48, finding=runs/pfas-rice-oos-multidataset: the K\_PL-gated lipid mechanism (fit on Yamazaki, NOT on Kim) predicts the independent K
  \item \texttt{evi-risk-grain-reliable} (experiment\_run): point=0.2, finding=runs/pfas-rice-oos-multidataset: on the reliable-detection Kim grain subset (DF>=15\%), the out-of-sample grain predictio
  \item \texttt{evi-risk-bound} (experiment\_run): point=0.7, finding=runs/pfas-rice-oos-lipid: even with the lipid mechanism, the worst-case out-of-sample grain miss is PFOS endosperm \textasciitilde{}5x u
\end{itemize}

\section{Discussion}
The verdict is a qualified yes, stated as an assurance LEVEL, not a binary. The model is usable as a SCREENING-level dietary risk-assessment tool: it now represents the full C4-C12 range with no structural blind spot (the breakthrough closes the long chains the single-pool core could not reach), and it predicts the dietary compartment out-of-sample within about a factor of 3 (reliable subset about 1.6) on an independent dataset using a mechanism that was not fit to that dataset. Two honesty conditions are first-class, not footnotes. (1) The structural coverage is a saturated reproduction (per-congener f\_xy and carrier are calibrated), so it certifies representability, not parameter-free prediction; in practice a risk assessor calibrates to a site, and the grain out-of-sample result is the predictive assurance. (2) The worst-case grain miss -- PFOS endosperm about 5x under, and GenX/ether over -- means the tool is screening-grade: dietary use must apply congener-specific uncertainty factors (about 3x typical, up to about 5x for PFSA endosperm and ether), and the long-chain carrier enhancement is not yet QSPR-able, so novel congeners inherit wider bounds. Within those bounds the model is fit for SCREENING dietary risk -- ranking congeners, flagging exceedances, prioritizing monitoring -- but not for regulatory-precision dose-response without further data (a reliable 2-pool root and an ether/sulfonamide QSPR are the named next requirements). The four claims reproduce under sci-adk verify; the breakthrough experiment is validation/longchain\_closure.py; the decision consolidates into FINDINGS.md.

\section{References}
No literature cited.

\end{document}